% Options for packages loaded elsewhere
\PassOptionsToPackage{unicode}{hyperref}
\PassOptionsToPackage{hyphens}{url}
\documentclass[
]{article}
\usepackage{xcolor}
\usepackage{amsmath,amssymb}
\setcounter{secnumdepth}{-\maxdimen} % remove section numbering
\usepackage{iftex}
\ifPDFTeX
  \usepackage[T1]{fontenc}
  \usepackage[utf8]{inputenc}
  \usepackage{textcomp} % provide euro and other symbols
\else % if luatex or xetex
  \usepackage{unicode-math} % this also loads fontspec
  \defaultfontfeatures{Scale=MatchLowercase}
  \defaultfontfeatures[\rmfamily]{Ligatures=TeX,Scale=1}
\fi
\usepackage{lmodern}
\ifPDFTeX\else
  % xetex/luatex font selection
\fi
% Use upquote if available, for straight quotes in verbatim environments
\IfFileExists{upquote.sty}{\usepackage{upquote}}{}
\IfFileExists{microtype.sty}{% use microtype if available
  \usepackage[]{microtype}
  \UseMicrotypeSet[protrusion]{basicmath} % disable protrusion for tt fonts
}{}
\makeatletter
\@ifundefined{KOMAClassName}{% if non-KOMA class
  \IfFileExists{parskip.sty}{%
    \usepackage{parskip}
  }{% else
    \setlength{\parindent}{0pt}
    \setlength{\parskip}{6pt plus 2pt minus 1pt}}
}{% if KOMA class
  \KOMAoptions{parskip=half}}
\makeatother
\usepackage{longtable,booktabs,array}
\newcounter{none} % for unnumbered tables
\usepackage{multirow}
\usepackage{calc} % for calculating minipage widths
% Correct order of tables after \paragraph or \subparagraph
\usepackage{etoolbox}
\makeatletter
\patchcmd\longtable{\par}{\if@noskipsec\mbox{}\fi\par}{}{}
\makeatother
% Allow footnotes in longtable head/foot
\IfFileExists{footnotehyper.sty}{\usepackage{footnotehyper}}{\usepackage{footnote}}
\makesavenoteenv{longtable}
\usepackage{graphicx}
\makeatletter
\newsavebox\pandoc@box
\newcommand*\pandocbounded[1]{% scales image to fit in text height/width
  \sbox\pandoc@box{#1}%
  \Gscale@div\@tempa{\textheight}{\dimexpr\ht\pandoc@box+\dp\pandoc@box\relax}%
  \Gscale@div\@tempb{\linewidth}{\wd\pandoc@box}%
  \ifdim\@tempb\p@<\@tempa\p@\let\@tempa\@tempb\fi% select the smaller of both
  \ifdim\@tempa\p@<\p@\scalebox{\@tempa}{\usebox\pandoc@box}%
  \else\usebox{\pandoc@box}%
  \fi%
}
% Set default figure placement to htbp
\def\fps@figure{htbp}
\makeatother
\usepackage{svg}
\ifLuaTeX
  \usepackage{luacolor}
  \usepackage[soul]{lua-ul}
\else
  \usepackage{soul}
\fi
\setlength{\emergencystretch}{3em} % prevent overfull lines
\providecommand{\tightlist}{%
  \setlength{\itemsep}{0pt}\setlength{\parskip}{0pt}}
\usepackage{bookmark}
\IfFileExists{xurl.sty}{\usepackage{xurl}}{} % add URL line breaks if available
\urlstyle{same}
\hypersetup{
  hidelinks,
  pdfcreator={LaTeX via pandoc}}

\author{}
\date{}

\begin{document}

\textbf{Behavioral finance tools for farmers' input price risk
management}

Riccardo Spada\textsuperscript{1}, Lukas Åkerman\textsuperscript{2},
Simone Cerroni\textsuperscript{3,4}, Miranda P.M.
Meuwissen\textsuperscript{1}, Laura Moritz\textsuperscript{3}, Jens
Rommel\textsuperscript{2}, Tobias Dalhaus\textsuperscript{1,5}

\textsuperscript{1} Business Economics Group, Wageningen University \&
Research, Netherlands

\textsuperscript{2} Department of Economics, Swedish University of
Agricultural Science (SLU), Sweden

\textsuperscript{3} Center Agriculture Food Environment (C3A),
University of Trento, Italy

\textsuperscript{4} Department of Economics and Management, University
of Trento, Italy

\textsuperscript{5} Agribusiness, Sustainable Food Systems Research
Centre, Rhine-Waal University of Applied Sciences, Kleve, Germany

\textbf{Abstract}

XXX

\textbf{JEL Codes:} XXX

\textbf{Keywords:} XXX

\textbf{Target Journals:}

\begin{itemize}
\item
  XXX
\end{itemize}

\section{Introduction}\label{introduction}

In recent decades, energy-intensive inputs in agriculture, such as
mineral fertilizers, have experienced considerable volatility, putting
farmers\textquotesingle{} overall income and, thereby, food security at
risk. Inorganic fertilizers, whose production requires substantial
amounts of energy, stand out as a critical component for modern
agriculture (European Commission, 2022; FAO, 2024). In particular,
nitrogen-based (N) fertilizers are the most important nutrient for crop
production (IFA, 2022). However, price shocks caused by recent
geopolitical crises are not completely transmitted to output prices
(Dong et al., 2025). Higher prices can reduce farmers' ability to apply
optimal nutrient levels, thereby lowering yields and jeopardizing food
security.

While in the long-run farms must reduce dependency on fossil energy
intensive inputs, price risk management can mitigate the financial
exposure to price volatility in the short run. A strategy to manage
price risk is to use financial tools such as forward contracts, futures
and options, and price insurance (Hardaker et al., 2015). But farmers'
adoption of such tools is lower than what would be expected from
utility-maximizing behavior of a risk-averse decision-maker(Spada et
al., 2025).

Many farmers' decisions deviate from pure profit maximization and are
better explained by considering risk preferences by using an expected
utility framework, such as fertilizer decision (Moritz et al., 2026).
Nonetheless, a growing body of empirical evidence shows that farmers'
choices often depart from the assumptions of Expected Utility Theory
(EUT). For example, producers display loss aversion and probability
weighting (Finger et al., 2024), consistent with Cumulative Prospect
Theory (CPT) (Tversky \& Kahneman, 1992). Additionally, adoption
decisions might be narrowly framed as stand-alone investments (Barberis
et al., 2006), instead of a risk management strategy within a broader
farm portfolio (broad framing). Therefore, CPT has been effectively used
to better explain adoption decisions of risk management tools (Mattos et
al., 2008; Babcock, 2015). Further, agricultural commodity price risk
management literature focuses mainly on output prices (Spada et al.,
2025), underscoring the need for further research on the input side.

In this paper, we aim to adjust contract specifics of financial tools
for farmers' input price risk management to better fit
farmers\textquotesingle{} behavioral preferences. Based on a simulation
exercise, we show how contract adjustments, such as introducing a price
ceiling, potentially affect farmers utility under EUT and CPT, as well
as under broad and narrow framing. This leads to a conclusion about
which contract specifics might be suitable under different farmers'
preferences and framing scenarios.

Our study extends the literature in two ways. First, there is a very
limited conceptualization of farmers' price risk management decisions
beyond standard expected utility theory (Spada et al., 2025). We thus
introduce a framework for input price risk management under risk using
different framings and preference scenarios, thereby paving the way for
experimental research to test for correlations between elicited
preferences (e.g. using XXX) and financial price risk management tool
adoption. Second, we contribute to the emerging literature on behavioral
contracts for farmers' risk management that better consider farmers'
actual decision making and preferences in contract design (e.g., Dalhaus
et al., 2020).

To achieve our objective, we first provide a theoretical background and
conceptualize how farmers might assess this decision under different
decision theories (EUT and CPT) and framings (broad and narrow). We then
build a simplified profit function of wheat production under risk.
Subsequently, we define contract specifications of price risk management
tools, that might affect farmers' evaluation of these tools in terms of
their utility under different behavioral models and framings. We then
employ a bio-economic model by integrating a crop simulation with
decision models under risk and ambiguity to test how adjusted contracts
affect utility values (Finger, 2012). First, we parametrize the
bio-economic model using simulated nitrogen and yield relationship from
a crop model and an econometric procedure to estimate stochastic
relations between nitrogen and yield. Second, we estimate nitrogen and
yield price dynamics. Lastly, we incorporate the obtained information in
the economic decision models and evaluate the effect on the outcome
distribution for the different models and framing.

We find that under standard EUT preferences the behaviorally adjusted
contract is attractive to farmers when broadly framed as a risk
management tool, but when the decision is narrowly framed as an
investment, we find that it doesn't provide any benefit. Under CPT
preferences, the adjusted contract is not attractive to farmers, both
when broadly and narrowly framed. We conclude that, while promising, to
make these types of contracts attractive to farmers, a more holistic
approach that considers also output prices might provide better results.
The remainder of the paper is organized as follows. In section 2, we lay
down the theoretical background. First, we review the most common
theories of farmers decision-making under risk and ambiguity. Second, we
develop a bio-economic production model for fertilizer price risk
decision-making. Building on this, we introduce different decision
framing and propose behavioral adjustments to risk management tools. In
section 3, we describe our methodology and in section 4, we present the
data and the assumptions used. We present the results of the
optimization model in section 5, and discuss them in section 6. We
conclude in section 7.

\section{Theoretical background}\label{theoretical-background}

\subsection{Theoretical frameworks on risk and
ambiguity}\label{theoretical-frameworks-on-risk-and-ambiguity}

Farmer decision-making behavior can be modeled using a variety of
theoretical frameworks that capture different attitudes toward risk and
ambiguity. We define risk as situations where the probabilities of
different outcomes are known, while ambiguity refers to situations where
these probabilities are not known (Knight, 1921). Traditional approaches
that model decisions under risk and ambiguity often rely on Expected
Utility Theory (EUT). However, empirical evidence suggests that
decision-makers may deviate from purely EUT maximizing behavior. To
capture these deviations, alternative models, most prominently
Cumulative Prospect Theory (CPT) by Tversky \& Kahneman (1992), have
been proven useful to better describe farmer behavior (Bocquého et al.,
2014; Babcock, 2015). CPT accounts for loss aversion, probability
weighting, and reference dependence. Unlike risk, where probabilities of
different outcomes can be estimated from historical data, ambiguity
occurs when the probabilities are not available or decision-makers do
not know them. In this context, the α-Expected Utility (αEU) model
(Ghirardato et al., 2004), by accounting for ambiguity aversion, might
be a better fit to explain farmers\textquotesingle{} behavior in
ambiguous situations.

Under standard EUT assumptions, a common approach is to model
farmers\textquotesingle{} behaviour as maximiser of certainty equivalent
(CE). The CE of a risky outcome \((x)\) denotes the guaranteed payoff
level that has the same utility as the expected utility of the risky
outcome. Risk averse farmers are willing to reduce risk exposure by
paying a risk premium. The expected outcome \(E(x)\) is thus reduced by
the risk premium as follows:

\[\begin{array}{r}
{CE = E(U) \approx}{\ E(x) - \ \frac{\delta}{2}Var(x)}\#(1)
\end{array}\]

where \(\delta\) is absolute risk aversion. While still relevant for
normative assessments, farmers decisions about financial tools are in
general not consistent with EUT (XXX).

Behavioral economics extensions of EUT might better explain actual
farmers behavior. Thus, assuming CPT preferences, where gains and losses
are considered as deviations from a reference point, decision-makers
might maximize the prospect value \((V)\) of a risky outcome as:

\[\begin{array}{r}
\max{V = \ }\sum_{i = 1}^{n}{\vartheta_{i}v\left( x_{i} \right)}\#(2)
\end{array}\]

where single prospect outcomes are transformed into a value function
\(v\left( x_{i} \right)\) that incorporates risk-aversion over gains and
risk-seeking behavior over losses:

\[\begin{array}{r}
v(x) = \left\{ \begin{array}{r}
{{(x}_{i} - x_{0})}^{a^{+}}\ \ \ \ \ \ \ \ \ \ \ \ \ \ \ \ \ \ \ ifx_{i} \geq x_{0} \\
 - \lambda\left( - \left( x_{i} - x_{0} \right) \right)^{a^{-}}\ \ \ \ \ ifx_{i} < x_{0}
\end{array} \right.\ \#(3)
\end{array}\]

where \(x_{0}\) is the reference point, \(\lambda\) is the loss-aversion
parameter, and \(a^{+}\) and \(a^{-}\) are the risk-aversion parameters
for gains and losses respectively. For each outcome, there is a
probability of occurrence \(\left( p_{i} \right)\). These probabilities
are translated into decision weights by a function \(w(p)\), which
captures the tendency to overweight small probabilities and underweight
large ones. Several functional forms have been proposed for the
probability weighting function, but the most commonly used are those
introduced by Tversky \& Kahneman (1992), which assumes different
weighting in the domains of gains and losses (Eq. 9), and Prelec (1998),
which assumes equal weighting in the two domains (Eq. 10):

\(\begin{array}{r}
w(p) = \left\{ \begin{array}{r}
\frac{p^{\gamma^{+}}}{\left\lbrack p^{\gamma^{+}} + (1 - p)^{\gamma^{+}} \right\rbrack^{1\text{/}\gamma^{+}\ }}\ \ \ \ if\ x \geq x_{0}\  \\
\frac{p^{\gamma^{-}}}{\left\lbrack p^{\gamma^{-}} + (1 - p)^{\gamma^{-}} \right\rbrack^{1\text{/}\gamma^{-}\ }}\ \ \ \ if\ x < x_{0}
\end{array} \right.\ \#(4)
\end{array}\)

\[\begin{array}{r}
w(p) = \exp\left\{ \  - \ \left\lbrack - \ln(p) \right\rbrack^{\gamma} \right\}\#(5)
\end{array}\]

in both cases, \(p\) is the objective probability and \(\gamma\) defines
the curvature of the function. Assuming ordered outcomes, the decision
weight for gains and losses are defined as:

\[\vartheta_{i}^{+} = w\left( p_{i} + \ldots + p_{n} \right) - w\left( p_{i + 1} + 1 + \ldots + p_{n} \right)\]

\[\begin{array}{r}
\vartheta_{i}^{-} = w\left( p_{m} + \ldots + p_{i} \right) - w\left( p_{m} + 1 + \ldots + p_{i - 1} \right)\#(6)
\end{array}\]

When decision-makers lack reliable information about the probabilities
of outcomes, decisions can occur under ambiguity. Several theories that
model decision-making under ambiguity have been proposed.\footnote{Conte
  \& Hey (2013) and Hey \& Pace (2014), assessed the most relevant
  decision models under ambiguity broadly dividing them into those that
  use second-order probabilities and those that do not. Among the
  non-SOMs, Hey \& Pace (2014) conclude that the Vector EU model
  (Siniscalchi, 2009) and the Alpha EU model (Ghirardato et al., 2004)
  perform better than the Subjective EU model (Savage, 1954) in
  explaining behaviour, but they perform similarly when it comes to
  prediction.} Here we adopt the αEU model proposed by Ghirardato et al.
(2004), that is a generalization of the MaxMin Expected Utility model
(Gilboa \& Schmeidler, 1989). These models are based on the assumption
that preferences can be represented by a set \(\mathcal{P}\) of beliefs
(i.e., probabilities) of different events, and that ambiguity can then
be described by the size of the belief set (Gilboa \& Schmeidler, 1989).
The decision-maker has then multiple subjective probabilities
\(P \in \mathcal{P}\) (or multiple priors). If the set \(\mathcal{P}\)
is a singleton, the decision-maker perceives no ambiguity and the model
reduces to the standard Subjective Expected Utility model. Moreover, the
\(\alpha\)EU model introduces a parameter \((\alpha)\) that capture the
degree of ambiguity attitudes. Decisions are then made on the basis of a
weighted average of the minimum (worst case) and the maximum (best case)
expected utility over the set \(\mathcal{P}\) of probabilities:

\[\begin{array}{r}
\max\ U = {\alpha\min_{P \in \mathcal{P}}{E^{P}U(x)} + (1 - \alpha)\ \max_{P \in \mathcal{P}}{E^{P}U(x)}}\#(7)
\end{array}\]

where \(\alpha \in \lbrack 0,1\rbrack\) is the decision-maker ambiguity
attitude. The larger \(\alpha\) is, the more the decision-maker is
ambiguity averse, and the larger the weight on the pessimistic
evaluation. In the special case where \(\alpha = 1\), the model reduces
to the MaxMin model of Gilboa \& Schmeidler (1989).

\subsection{\texorpdfstring{Decision framing
}{Decision framing }}\label{decision-framing}

Individuals can evaluate risky decisions using different mental framing
(Tversky \& Kahneman, 1981). With a broad framing, decision-makers might
consider a risky choice in conjunction with other risks. In this case
input contracting is considered as a tool for managing farming risk, and
utility is defined over total wealth or profits. In contrast, under
narrow framing (Barberis et al., 2006), decision-makers might consider
the same choice independently from other risks. In this case input
contracting is viewed as a stand-alone investment and utility is defined
by taking into account each choice in isolation, rather than part of a
broader farm strategy. Although narrow framing might lead to biased
decisions, it might better explain farmers behaviour. For example it has
been effective in explaining observed crop insurance choices (Babcock,
2015).

\subsection{\texorpdfstring{Economic model
}{Economic model }}\label{economic-model}

To conceptualize farmers' input price risk management decision we build
a simplified model of crop production for winter wheat, with (synthetic)
N fertilizer as the single input, and the associated risks and
uncertainties. Following the initial crop choice, farmers make
fertilizer purchasing decisions under risk (Figure 1), given the price
fluctuations until the fertilizer is used. By committing to buy a
specific quantity of fertilizer at an agreed price and future date,
forward contracting effectively mitigates exposure to price risk
(Hardaker et al., 2015).\footnote{Other contract specifications may also
  include quality and time of delivery, and input suppliers can also
  offer additional benefits such as financing and fertiliser
  recommendations.} However, this decision might not be only influenced
by fertilizer price risk. When considering fertilizer-related decisions,
farmers might take into account the impact of this choice on total
profits (broad framing). Here, other sources of uncertainty enter the
decision problem such as yield, and output prices, making the decision
ambiguous.

\begin{figure}
\centering
\includesvg[width=6.5in,height=1.71875in]{media/image2.svg}
\caption{\emph{Figure 1. Economic decision framework}}
\end{figure}

We consider a static model of a representative farmer who produces a
single commodity (winter wheat) using a single input (synthetic N
fertilizer). Expected farmer profits \((\pi)\) are defined as revenue
minus costs:

\[\begin{array}{r}
\pi = p_{w}y(N) - p_{N}N\#(8)
\end{array}\]

where \(p_{w}\) is the wheat price, \(y\) the expected yield, \(p_{N}\)
the fertilizer price, and \(N\) the fertilizer quantity. We assume that,
once the crop is determined at time \(t_{0}\), the farmer knows the
associated fertilizer needs for the season. All the volume is then
either purchased at market price or through a forward contract. When
farmers simply buy fertilizers close to application, they pay the
stochastic spot price \(\left( {\widetilde{p}}_{N,1} \right)\). When
forward contracting, \footnote{We do not specify other contract
  parameters like quantity, delivery, and settlement date, as they can
  be flexibly arranged for both cash and forward purchases.} the input
costs become fixed \((f)\), with the forward price calculated as:

\(\begin{array}{r}
f = p_{N,0}e^{- rT}\#(9)
\end{array}\)where \(p_{N,0}\) is the fertilizer spot price at
\(t_{0}\), \(r\) is the risk-free rate, and \(T\) is the delivery time.
We also assume no fixed and variable costs (e.g., information gathering
expenses or commissions) for entering the contract. We consider forward
contracting as mathematically equivalent to a buy-and-hold strategy. If
the input is bought at time \(t_{0}\), the cost at time \(t_{1}\) is the
current spot price plus storage costs. With a forward contract, the firm
instead commits to a purchase price \((f)\) that already reflects these
same costs. Hence, forward contracting reproduces the same outcome as
purchasing and store, while avoiding holding inventory. In this context
we assume the cost of forward contracting for fertilizers to be zero.

\subsection{Behaviorally adjusted contract
design}\label{behaviorally-adjusted-contract-design}

In addition to the option of buying fertilizers at the spot price
\(\left( {\widetilde{p}}_{N,1} \right)\) and forward contracting with a
supplier, thus fixing the price \((f)\), we introduce an ``insurance''
mechanism to the contract that sets a price ceiling. By paying a premium
\((\Gamma)\), if the fertilizer price is above a certain threshold
\((k)\), the insurance payoff \((I)\) would ensure that the price
farmers pay for fertilizer would be capped. This mechanism thus
functions as a European call option where the payout amount is made
whenever the fertiliser spot price at time \(t = 1\) is above the
farm\textquotesingle s individual threshold price, otherwise the payout
is zero:

\[\begin{array}{r}
I = \max{(p_{N,1} - k,0)}\ \ \#(10)
\end{array}\]

In this case, and unlike the traditional forward contract, the contract
is not free to enter. Here, the supplier carries all the risk and should
thus be compensated. To calculate the premium of an European call
option, we discount its expected payoff:

\[\begin{array}{r}
\Gamma = E(I)e^{- rT}\#(11)
\end{array}\]

\protect\phantomsection\label{_Ref205816026}{}Traditional forward
contracts primarily address risk aversion by fixing the price and
eliminating exposure to future price volatility. In contrast, the
proposed adjusted contract (also called maximum price contract) targets
downside risk aversion and loss aversion by offering protection against
adverse price movements while preserving the opportunity to benefit from
price declines. Available purchasing choices are summarized in
\hyperref[_Ref205816026]{Table 1}.

\begin{longtable}[]{@{}
  >{\raggedright\arraybackslash}p{(\linewidth - 6\tabcolsep) * \real{0.2490}}
  >{\centering\arraybackslash}p{(\linewidth - 6\tabcolsep) * \real{0.1601}}
  >{\centering\arraybackslash}p{(\linewidth - 6\tabcolsep) * \real{0.2120}}
  >{\centering\arraybackslash}p{(\linewidth - 6\tabcolsep) * \real{0.3788}}@{}}
\caption{\textbf{Table 1.} \emph{Overview of traditional and adjusted
contract design}}\tabularnewline
\toprule\noalign{}
\begin{minipage}[b]{\linewidth}\raggedright
\end{minipage} & \begin{minipage}[b]{\linewidth}\centering
\textbf{Spot price}
\end{minipage} & \begin{minipage}[b]{\linewidth}\centering
\textbf{Forward contract}
\end{minipage} & \begin{minipage}[b]{\linewidth}\centering
\textbf{Adjusted contract}
\end{minipage} \\
\midrule\noalign{}
\endfirsthead
\toprule\noalign{}
\begin{minipage}[b]{\linewidth}\raggedright
\end{minipage} & \begin{minipage}[b]{\linewidth}\centering
\textbf{Spot price}
\end{minipage} & \begin{minipage}[b]{\linewidth}\centering
\textbf{Forward contract}
\end{minipage} & \begin{minipage}[b]{\linewidth}\centering
\textbf{Adjusted contract}
\end{minipage} \\
\midrule\noalign{}
\endhead
\bottomrule\noalign{}
\endlastfoot
Fertilizer cost & \({\widetilde{p}}_{N,1}\) & \(p_{N,0}e^{- rT}\) &
\(\left\{ \begin{array}{r}
p_{N,1} - k\ \ \ if\ p_{N,1} > k \\
p_{N,1}\ \ \ \ \ \ \ \ \ \ \ if{\ p}_{N,1} \leq k
\end{array} \right.\ \) \\
Contract cost & \emph{-} & \emph{-} & \(\Gamma\) \\
Mechanism & \emph{-} & \emph{Price fixing} & \emph{Price ceiling} \\
Behavior targeted & \emph{-} & \emph{Risk aversion} & \emph{Downside
risk aversion/Loss aversion} \\
\end{longtable}

Notes: we assume no cost of storage for the forward contract.

\subsection{Profit functions}\label{profit-functions}

In our model farmers can choose whether to buy part or all their
fertilizers needs through a forward contract or the adjusted contract.
For a forward contract, farmers perceive the value of their contract
choice as positive or negative (i.e. gain or loss) depending on whether
they paid more or less than the underlying spot price. In this case, the
payoff is calculated as:

\[\begin{array}{r}
\Pi_{F} = f - p_{N}\#(12)
\end{array}\]

For the adjusted contract, a gain (loss) occurs when the payoff received
from the insurance mechanism \((I)\) is greater (smaller) than the
premium paid \((\Gamma)\). The adjusted contract benefits are then
calculated as the payoff minus the premium paid:

\[\begin{array}{r}
\Pi_{I} = I - \Gamma\#(13)
\end{array}\]

In the case of broad framing, profits are then defined as:

\[\begin{array}{r}
\pi_{BF} = p_{w}y(N) - p_{N}N + \Pi_{i}N\#(14)
\end{array}\]

with \(i = \lbrack 0,F,I\rbrack\), in the case of no contract, forwards,
or adjusted contract respectively. In the case of narrow framing,
utility is narrowly defined by considering only the decision at hand:

\[\begin{array}{r}
\pi_{NF} = \Pi_{i}N\#(15)
\end{array}\]

\section{Methodology}\label{methodology}

In this section we present the different economic decision models, the
biophysical and statistical approaches used to estimate the models, and
the computation method of the adjusted contract price. For the decision
under broad framing we first estimate the interactions between nitrogen
application and crop yield, and the relationship between nitrogen and
wheat prices. The data sources and the estimation will be discussed in
the subsequent sections.

\subsection{\texorpdfstring{Nitrogen -- yield relationship
}{Nitrogen -- yield relationship }}\label{nitrogen-yield-relationship}

We estimate nitrogen-yield relationships using data from the
Environmental Policy Integrated Climate-based (EPIC; Williams 1995) crop
model, where different level of nitrogen are used to simulate crop
yield. Simulated observations of winter wheat yield for different levels
of nitrogen applications are then used in a production function to
estimate the statistical relationship between nitrogen and yield under
risk. The estimation employed here, based on Just \& Pope (1978, 1979),
allows inputs to influence both the mean and the variability of outputs:

\[\begin{array}{r}
y(N) = f(N) + \sigma_{y}(N)\varepsilon\#(16)
\end{array}\]

where \(f(N)\) represents the expected yield (production function) and
\(\sigma_{y}(N)\) the standard deviation of yield (variation function),
with the assumption that \(E(\varepsilon) = 0\) and
\(\sigma(\varepsilon) = 1\). Following Finger \& Hediger
(2008),\footnote{Finger \& Hediger (2008) analyzed different functional
  form for Swiss corn production and suggest that a square root
  specification of the production function leads to the smallest cost of
  misspecification.} we employ a production function estimation of the
form:

\[\begin{array}{r}
f(N) = \beta_{0} + \beta_{1}N^{0.5} + \beta_{2}N + e\#(17)
\end{array}\]

If \(\beta_{1} > 0\) and \(\beta_{2} < 0\), the function shows
decreasing marginal productivity of nitrogen, where yields are
monotonically increasing up to some point of nitrogen use and then
monotonically decreasing. We then define yield variation as the absolute
difference between observed and expected yields (Finger \& Schmid,
2008).\footnote{Observed yields are yields simulated with the crop
  model, expected yields are derived from the production function.}
Therefore, we regress residual from the production function estimation
in (17) (\(\widehat{e} = y - \widehat{y})\) to estimate the yield
variation function, i.e. the relation between the level of nitrogen
input and yield variance:

\[\begin{array}{r}
\sigma_{y}(N) = \left| \widehat{e} \right| = \beta_{0} + \beta_{1}N^{0.5}\#(18)
\end{array}\]

The production and yield variability function are estimated with the
MM-estimator (Yohai, 1987) to reduce potential influences of outlying
observations on estimation results (Finger, 2013).

\subsection{\texorpdfstring{Nitrogen-Yield prices relationship
}{Nitrogen-Yield prices relationship }}\label{nitrogen-yield-prices-relationship}

To estimate the joint dynamic price processes of nitrogen \((p_{N})\)
and wheat \((p_{w})\), we use a quantile vector autoregression (QVAR)
and a copula approach (Li \& Chavas, 2023). Stochastic price dynamics
are represented by a general autoregressive process of the form:

\[\begin{array}{r}
\mathbf{p}_{t} = \begin{bmatrix}
p_{N,t} \\
p_{w,t}
\end{bmatrix} = \begin{bmatrix}
g_{N}(\mathbf{p}_{t - 1},\ldots,\mathbf{p}_{t - n} \\
g_{w}(\mathbf{p}_{t - 1},\ldots,\mathbf{p}_{t - n}
\end{bmatrix}\#(19)
\end{array}\]

where the distribution functions 𝑔(⋅) explain the contemporaneous prices
through their lags.\footnote{With the assumption that 𝑔(⋅) is
  differentiable, and its error is serially independent with a
  continuous distribution function.} We are interested in the joint
distribution function described as:

\[\begin{array}{r}
F\left( \mathbf{p} \middle| \mathbf{P}_{t - 1} \right) = Prob\lbrack\mathbf{g}(\mathbf{p}_{t - 1},\ldots,\mathbf{p}_{t - n} \leq \mathbf{p}\rbrack\#(20)
\end{array}\]

where
\(\mathbf{P}_{t - 1} = (\mathbf{p}_{t - 1},\ldots,\mathbf{p}_{t - n})\)
includes lagged realizations of all prices. To estimate the joint
distribution \(F\), we first derive the marginal distribution functions
using a QVAR model:

\[\begin{array}{r}
F_{N}\left( p_{N} \middle| \mathbf{P}_{t - 1} \right) = Prob\lbrack\mathbf{g}_{N}(\mathbf{p}_{t - 1},\ldots,\mathbf{p}_{t - n} \leq p_{N}\rbrack\#(21a)
\end{array}\]

\[\begin{array}{r}
F_{w}\left( p_{w} \middle| \mathbf{P}_{t - 1} \right) = Prob\lbrack\mathbf{g}_{w}(\mathbf{p}_{t - 1},\ldots,\mathbf{p}_{t - n} \leq p_{w}\rbrack\#(21b)
\end{array}\]

We then rely on Sklar's theorem (Sklar, 1959), and use a copula \(C\)
(Nelsen, 2006) to estimate the following relationship:

\[\begin{array}{r}
F\left( \mathbf{p} \middle| \mathbf{P}_{t - 1} \right) = C\left\lbrack F_{N}\left( p_{N} \middle| \mathbf{P}_{t - 1} \right),F_{w}\left( p_{w} \middle| \mathbf{P}_{t - 1} \right) \right\rbrack\#(22)
\end{array}\]

\subsection{\texorpdfstring{Adjusted contract pricing
}{Adjusted contract pricing }}\label{adjusted-contract-pricing}

We derive the contract pricing with a burn analysis. We implement this
by determining the discounted average payoff of the contract over \(n\)
simulations. The general pricing formula (Eq. 11) is then implemented as
follows:

\[\begin{array}{r}
\Gamma = \frac{1}{n}\sum_{t - 1}^{n}(I)e^{- rT}\#(23)
\end{array}\]

We specify the assumptions needed for the calculation following our
decision framework (section \hl{}\hyperref[economic-model]{2.3}). The
planning horizon is set for one season, while the considered contract
duration is six months (from crop choice/planting in early October to
the first fertilizer application in early March). Different strike
prices will be analyzed based on the price distribution, allowing for
evaluation of contracts across different price protection threshold. A
standard risk free interest rate of 3\% is applied to discount future
payoffs to present values.

\subsection{Montecarlo simulation of utility under different decision
theories and
framing}\label{montecarlo-simulation-of-utility-under-different-decision-theories-and-framing}

The previously described information on yield, price dynamics, and
contract cost, enter the different economic models of farmers'
decision-making. More specifically, we analyse the impact on utility
level for different decision theories and different framing. In the case
farmers rely on a broad decision framing, Eq. (12) will be evaluated
under EUT, CPT, and the α-EU frameworks. Under narrow framing, farmers
profit are instead derived from Eq. (15). Given that no ambiguity arise
under narrow framing, only EUT and CPT preferences are considered.

Assumptions regarding the representative farm characteristics, and the
behavioural parameters entering the models have to be specified. For a
German representative arable farm, average N fertilizer requirement are
200 kgN/ha, sowing is usually in late October, and harvest in early
August (Hernández-Ochoa et al., 2023). Absolute risk aversion for the
EUT specification \((\delta)\) is considered between 0.00001 (less
risk-averse) and 0.001 (more risk-averse). Structural parameter for CPT
and \(\alpha\)EU are derived from experimental evidence available in the
literature. For the CPT model, final prospect values are dependent on
risk aversion \((a)\), loss aversion \((\lambda)\), and probability
distortion \((\gamma)\). Broad evidence on structural parameter
elicitation is available for European farmers in general and German
farmers in particular (Bougherara et al., 2017; Finger et al., 2024;
Rommel et al., 2023). The considered parameters values are:
\(a = 0.334\), \(\lambda = 1.574\), and \(\gamma = 0.571\). Alternative
parametrization are considered for robustness checks. As reference
points, we follow Kőszegi \& Rabin (2007) argument to use expected
outcomes in the case of broad framing (in our case expected profit per
hectare), and zero for narrow framing. For the \(\alpha\)EU, we consider
a binary state space of either high or low profits occurring with
probability \(P\). Ambiguity arises as the farmer knows that \(P\) lies
in a certain interval
\(\mathcal{P =}\left\lbrack P_{\min},P_{\max} \right\rbrack\) that we
naively derive from profits distribution. The parameter for ambiguity
aversion \((\alpha)\) is considered to be
\(\lbrack 0.1,\ 0.5,\ 0.9\rbrack\), reflecting ambiguity loving,
neutrality, and aversion respectively (Couture et al., 2024).

\section{Data}\label{data}

In this section we present: i) the dataset used to estimate the
nitrogen-yield relationship, and ii) the price data. We obtained
fertilizers and wheat prices farm level data for the Cologne Lowland
area, in the southwest of the German Federal state of North
Rhine-Westphalia. We thus focus our analysis of fertilizer-yield
relationship on the same geographic area.

\subsection{\texorpdfstring{Yield data }{Yield data }}\label{yield-data}

We use data generated from the European implementation of the EPIC model
(EPIC-IIASA; Balkovič et al., 2013). We focus on a representative
location in the Cologne Lowland area, a highly productive region
specialized on crop farms. The dataset includes winter wheat yields for
different levels of N applications, management practices and weather
scenarios. \hyperref[_Ref212044160]{Table 2} presents summary statistics
of generated yield levels for different N application rates. It shows a
linear relationship between N application rates and the level of wheat
yield.

Table A1 in the appendix shows complete descriptive statistics for the
simulated data. Different water and land management practices as well as
different weather scenarios are included in the simulation. Irrigation
is applied for N application rates \(\geq\) 180 kg N/ha. In the case of
monoculture, conventional tillage with different percentages of crop
residues left in the field are assumed across 0, 30, 60, and 90 percent.
In the case of different crop rotations, conventional tillage with
business-as-usual crop residue and minimum tillage with high crop
residues are considered. Original dry-matter yield levels are converted
to economic yield considering a moisture content of 12\%. Yields for
different N levels are generated for the period 1985-2020 based on
historical weather. To increase the number of scenarios, we moreover
consider simulated yields under future climate scenarios for different
emission pathways. More specifically, we consider a low emission
(ssp126) and a very high emission scenario (ssp585) for the simulation
of yields only for N application rates of 180 kgN/ha.\footnote{Based on
  phase 6 of the Coupled Model Intercomparison Project (cmip6). Outputs
  from the Institut Pierre-Simone Laplace (ipsl) and Max Planck
  Institute for Meteorology (mpi) climate models are considered to
  simulate yields under different Shared Socioeconomic Pathways (ssp);
  ssp126 is a low-emission scenario; ssp585 is a very high emission
  scenario (O'Neill et al., 2014).} For more information about the
EPIC-IIASA model specifications and calibration see Balkovič et al.
(2013).

\protect\phantomsection\label{_Ref212044160}{}\emph{\hfill\break
}

\emph{Table 2.} \emph{Summary statistics of simulated wheat yield
(t//ha) from the EPIC-IIASA crop model}

{\def\LTcaptype{none} % do not increment counter
\begin{longtable}[]{@{}
  >{\raggedright\arraybackslash}p{(\linewidth - 10\tabcolsep) * \real{0.1892}}
  >{\centering\arraybackslash}p{(\linewidth - 10\tabcolsep) * \real{0.0937}}
  >{\centering\arraybackslash}p{(\linewidth - 10\tabcolsep) * \real{0.2034}}
  >{\centering\arraybackslash}p{(\linewidth - 10\tabcolsep) * \real{0.0614}}
  >{\centering\arraybackslash}p{(\linewidth - 10\tabcolsep) * \real{0.0715}}
  >{\centering\arraybackslash}p{(\linewidth - 10\tabcolsep) * \real{0.1450}}@{}}
\toprule\noalign{}
\begin{minipage}[b]{\linewidth}\raggedright
\textbf{N application rate}

\textbf{(kgN/ha)}
\end{minipage} & \begin{minipage}[b]{\linewidth}\centering
\textbf{Median}
\end{minipage} & \begin{minipage}[b]{\linewidth}\centering
\textbf{Median of Absolute}

\textbf{Deviation (MAD)\textsuperscript{a}}
\end{minipage} & \begin{minipage}[b]{\linewidth}\centering
\textbf{Min}
\end{minipage} & \begin{minipage}[b]{\linewidth}\centering
\textbf{Max}
\end{minipage} & \begin{minipage}[b]{\linewidth}\centering
\textbf{Observations}
\end{minipage} \\
\midrule\noalign{}
\endhead
\bottomrule\noalign{}
\endlastfoot
10 & 1.6 & 0.18 & 1.16 & 2.31 & 35 \\
50 & 3.65 & 0.69 & 2.02 & 4.65 & 140 \\
100 & 5.76 & 0.87 & 3.05 & 7.15 & 140 \\
180 & 7.83 & 1.01 & 4.30 & 14.01 & 1084 \\
250 & 11.23 & 0.86 & 4.68 & 13.32 & 210 \\
\end{longtable}
}

\textsuperscript{a} scaled by a correction factor of 1.4826

\subsection{Price data}\label{price-data}

We use farm-level fertilizer and wheat prices data obtained from the
Chamber of Agriculture of North Rhine-Westphalia. We use Calcium
Ammonium Nitrate (CAN) as representative N-fertilizer that is widely
applied and for which data is most complete. Prices for CAN are recorded
twice a month (1\textsuperscript{st} and 3\textsuperscript{rd} Monday of
the month)\footnote{This sampling frequency generates calendar-related
  time gaps of 21 days (or 3 weeks). For analysis purposes, we assume
  this interval is consistent across years and does not affect the
  interpretation of temporal trends.}, while the wheat price time series
contains weekly observations. We thus only use wheat price observations
in weeks for which we have also fertilizer price recordings. The final
dataset consists of 401 biweekly observations of CAN and wheat prices
from January 2009 to November 2025 (\hyperref[_Ref213337244]{Figure 2}).
For our analysis we use the first difference as both CAN and wheat time
series are not stationary\footnote{To check for stationarity we run
  three different tests: the Augmented Dickey-Fuller (ADF) test, the
  Kwiatkowski-Phillips-Schmidt-Shin (KPSS), and the Phillips-Perron (PP)
  test.} (see Table A2 in appendix).

\begin{figure}
\centering
\includegraphics[width=6.5in,height=5.17708in]{media/image3.png}
\caption{\protect\phantomsection\label{_Ref213337244}{}\emph{Figure 2.
Calcium Ammonium Nitrate (CAN) and wheat log of prices (€/kg; Jan
2009-Nov 2025).}}
\end{figure}

\section{Results}\label{results}

\subsection{Production and yield variation
functions}\label{production-and-yield-variation-functions}

The estimations of the production and yield variation functions are
presented in \hyperref[_Ref214287463]{Table 3}. Variables to control for
land and water management are included. Coefficient estimates of the
production function indicate that N application does not display
diminishing marginal productivity (\(\beta_{N} > 0\) within the observed
application range) and might be due to the model simulating only optimal
nitrogen application levels. The variation function shows that yield
variability (i.e. production risk) increases with higher nitrogen
applications.

\protect\phantomsection\label{_Ref214287463}{}\emph{Table 3.}
\emph{Estimation of production (Eqs. 17) and variation (Eqs. 18)
functions}

{\def\LTcaptype{none} % do not increment counter
\begin{longtable}[]{@{}
  >{\raggedright\arraybackslash}p{(\linewidth - 4\tabcolsep) * \real{0.2672}}
  >{\raggedright\arraybackslash}p{(\linewidth - 4\tabcolsep) * \real{0.3794}}
  >{\raggedright\arraybackslash}p{(\linewidth - 4\tabcolsep) * \real{0.3534}}@{}}
\toprule\noalign{}
\begin{minipage}[b]{\linewidth}\raggedright
\end{minipage} & \begin{minipage}[b]{\linewidth}\raggedright
\textbf{Production function}
\end{minipage} & \begin{minipage}[b]{\linewidth}\raggedright
\textbf{Variation function}
\end{minipage} \\
\midrule\noalign{}
\endhead
\bottomrule\noalign{}
\endlastfoot
Intercept & -0.464** & 0.087 \\
& (0.205) & (0.116) \\
N\textsuperscript{0.5} & 0.284*** & 0.081*** \\
& (0.049) & (0.025) \\
N & 0.026*** & -0.003** \\
& (0.003) & (0.001) \\
irrigation & 0.118* & -0.053 \\
& (0.069) & (0.039) \\
til\_min\_cons & -1.596*** & -0.040 \\
& (0.072) & (0.044) \\
til\_contill\_r00 & 1.013*** & -0.243*** \\
& (0.061) & (0.035) \\
til\_contill\_r30 & 0.632*** & -0.042 \\
& (0.093) & (0.051) \\
til\_contill\_r60 & 0.407** & 0.226*** \\
& (0.167) & (0.062) \\
til\_contill\_r90 & 0.030 & 0.355*** \\
& (0.196) & (0.075) \\
Num.Obs. & 1609 & 1609 \\
R2 & 0.864 & 0.098 \\
R2 Adj. & 0.863 & 0.093 \\
\end{longtable}
}

Note: til\_min\_cons: dummy indicating minimum tillage with high crop
residues left in the field; til\_contill\_r: dummies indicating
conventional tillage with different level of crop residues left in the
field (0, 30, 60, and 90\%); statistics in parentheses are standard
errors; significance at * = 0.1, ** = 0.05, and *** = 0.01 levels.

\subsection{Marginal price
distributions}\label{marginal-price-distributions}

We use QVAR to estimate the marginal distributions (21 a-b). To
determine the number of lags, we use the Schwarz Bayesian Information
Criterion (BIC). As suggested by the BIC, a model of order 1 was
selected for parsimony. The QVAR specification allows price dependencies
to be different across price quantiles. The estimation is thus carried
out for each integer quantile \(q\) from 1 to 99 (Koenker, 2005), with
the following specification:

\[\begin{array}{r}
p_{N}\left( q_{N} \middle| \mathbf{P}_{t - 1} \right) = \mathbf{\beta}_{\mathbf{0,\ \ }N}\left( q_{N} \right) + \mathbf{\beta}_{\mathbf{1,\ \ }N}\left( q_{N} \right)\mathbf{P}_{t - 1}\#(24a)
\end{array}\]

\[\begin{array}{r}
p_{w}\left( q_{w} \middle| \mathbf{P}_{t - 1} \right) = \mathbf{\beta}_{\mathbf{0,\ \ }w}\left( q_{w} \right) + \mathbf{\beta}_{\mathbf{1,\ \ }w}\left( q_{w} \right)\mathbf{P}_{t - 1}\#(24b)
\end{array}\]

where \(\mathbf{\beta}_{\mathbf{0,\ \ }i}(q_{i})\), with
\(i \in \{ N,w\}\), describes the intercept for price \(p_{i}\) in
quantile \(q_{i}\) and
\(\mathbf{\beta}_{\mathbf{1,\ \ }i}(q_{i})\mathbf{P}_{t - 1}\) is a
matrix of parameters for the conditional effect of each lag and
quantile. We investigate various functional specifications of the QVAR
mode by comparing information criteria on the mean regression
(VAR).\footnote{Exogenous variables investigated are squared endogenous
  variables, indexed time trends, and seasonal dummies. Information
  criteria improve with the addition of squared terms but we proceed
  with the linear model, as the non-linear specification does not
  produce a stable price simulation.} Results of the estimation of the
VAR and QVAR for selected quantiles (\(q\) = 0.1, 0.3, 0.5, 0.7, 0.9)
are reported in \hyperref[_Ref214526676]{Table 4} and
\hyperref[_Ref214526730]{Table 5}. Standard errors are derived by
bootstrapping (Koenker \& Xiao, 2006), and the goodness of fit for the
QVAR is calculated following Koenker \& Machado (1999).

\protect\phantomsection\label{_Ref214526676}{}\emph{Table 4.} \emph{VAR
and QVAR estimates for Calcium Ammonium Nitrate price}

{\def\LTcaptype{none} % do not increment counter
\begin{longtable}[]{@{}
  >{\raggedright\arraybackslash}p{(\linewidth - 12\tabcolsep) * \real{0.1528}}
  >{\raggedright\arraybackslash}p{(\linewidth - 12\tabcolsep) * \real{0.1412}}
  >{\raggedright\arraybackslash}p{(\linewidth - 12\tabcolsep) * \real{0.1412}}
  >{\raggedright\arraybackslash}p{(\linewidth - 12\tabcolsep) * \real{0.1412}}
  >{\raggedright\arraybackslash}p{(\linewidth - 12\tabcolsep) * \real{0.1412}}
  >{\raggedright\arraybackslash}p{(\linewidth - 12\tabcolsep) * \real{0.1412}}
  >{\raggedright\arraybackslash}p{(\linewidth - 12\tabcolsep) * \real{0.1410}}@{}}
\toprule\noalign{}
\begin{minipage}[b]{\linewidth}\raggedright
\end{minipage} & \begin{minipage}[b]{\linewidth}\raggedright
\end{minipage} & \begin{minipage}[b]{\linewidth}\raggedright
\end{minipage} & \begin{minipage}[b]{\linewidth}\raggedright
\[p_{N}\]
\end{minipage} & \begin{minipage}[b]{\linewidth}\raggedright
\end{minipage} & \begin{minipage}[b]{\linewidth}\raggedright
\end{minipage} & \begin{minipage}[b]{\linewidth}\raggedright
\end{minipage} \\
\midrule\noalign{}
\endhead
\bottomrule\noalign{}
\endlastfoot
& VAR & & & QVAR & & \\
Quantiles & & 0.1 & 0.3 & 0.5 & 0.7 & 0.9 \\
Intercept & -0.000 & -0.096*** & -0.026*** & 0.000 & 0.023*** &
0.091*** \\
& (0.011) & (0.016) & (0.004) & (0.003) & (0.005) & (0.011) \\
\(p_{N,t - 1}\) & 0.252*** & 0.362*** & 0.328*** & 0.312*** & 0.303*** &
0.274*** \\
& (0.047) & (0.089) & (0.050) & (0.030) & (0.053) & (0.026) \\
\(p_{w,t - 1}\) & 0.485*** & 0.056 & 0.086*** & 0.056*** & 0.099** &
0.223* \\
& (0.094) & (0.083) & (0.026) & (0.010) & (0.038) & (0.130) \\
Num.Obs. & 396 & 396 & 396 & 396 & 396 & 396 \\
Adj. R\textsuperscript{2} & 0.133 & & & & & \\
Pseudo R\textsuperscript{2} & & 0.089 & 0.094 & 0.064 & 0.073 & 0.084 \\
RMSE & 0.21 & 0.24 & 0.22 & 0.22 & 0.22 & 0.23 \\
\end{longtable}
}

Note: quantile regression standard errors are calculated by
bootstrapping; stars indicate significance at at * = 0.1, ** = 0.05, and
*** = 0.01 levels; Pseudo R\textsuperscript{2} for quantile regression
are calculated following Koenker \& Machado (1999).

\protect\phantomsection\label{_Ref214526730}{}\emph{Table 5.} \emph{VAR
and QVAR estimates for wheat price}

{\def\LTcaptype{none} % do not increment counter
\begin{longtable}[]{@{}
  >{\raggedright\arraybackslash}p{(\linewidth - 12\tabcolsep) * \real{0.1486}}
  >{\raggedright\arraybackslash}p{(\linewidth - 12\tabcolsep) * \real{0.1380}}
  >{\raggedright\arraybackslash}p{(\linewidth - 12\tabcolsep) * \real{0.1476}}
  >{\raggedright\arraybackslash}p{(\linewidth - 12\tabcolsep) * \real{0.1476}}
  >{\raggedright\arraybackslash}p{(\linewidth - 12\tabcolsep) * \real{0.1420}}
  >{\raggedright\arraybackslash}p{(\linewidth - 12\tabcolsep) * \real{0.1380}}
  >{\raggedright\arraybackslash}p{(\linewidth - 12\tabcolsep) * \real{0.1380}}@{}}
\toprule\noalign{}
\begin{minipage}[b]{\linewidth}\raggedright
\end{minipage} & \begin{minipage}[b]{\linewidth}\raggedright
\end{minipage} & \begin{minipage}[b]{\linewidth}\raggedright
\end{minipage} & \begin{minipage}[b]{\linewidth}\raggedright
\[p_{w}\]
\end{minipage} & \begin{minipage}[b]{\linewidth}\raggedright
\end{minipage} & \begin{minipage}[b]{\linewidth}\raggedright
\end{minipage} & \begin{minipage}[b]{\linewidth}\raggedright
\end{minipage} \\
\midrule\noalign{}
\endhead
\bottomrule\noalign{}
\endlastfoot
& VAR & & & QVAR & & \\
Quantiles & & 0.1 & 0.3 & 0.5 & 0.7 & 0.9 \\
Intercept & 0.001 & -0.086*** & -0.028*** & 0.000 & 0.029*** &
0.094*** \\
& (0.006) & (0.009) & (0.005) & (0.003) & (0.004) & (0.009) \\
\(p_{N,t - 1}\) & 0.075*** & 0.066 & 0.077*** & 0.086** & 0.050** &
0.053** \\
& (0.025) & (0.056) & (0.018) & (0.038) & (0.024) & (0.026) \\
\(p_{w,t - 1}\) & -0.119** & -0.037 & 0.099 & 0.099*** & 0.037 &
-0.060 \\
& (0.050) & (0.069) & (0.063) & (0.015) & (0.064) & (0.128) \\
Num.Obs. & 396 & 396 & 396 & 396 & 396 & 396 \\
Adj. R\textsuperscript{2} & 0.027 & & & & & \\
Pseudo R\textsuperscript{2} & & 0.032 & 0.029 & 0.024 & 0.015 & 0.016 \\
RMSE & 0.11 & 0.14 & 0.12 & 0.12 & 0.12 & 0.15 \\
\end{longtable}
}

Note: quantile regression standard errors are calculated by
bootstrapping; stars indicate significance at * = 0.1, ** = 0.05, and
*** = 0.01 levels; Pseudo R\textsuperscript{2} for quantile regression
are calculated following Koenker \& Machado (1999).

Results indicate that price changes for CAN are correlated to the lag of
own price changes and mostly correlated to the lag of wheat price
changes, but more around the central quantiles. Price changes for wheat
are correlated to own lagged price changes only at the median, while
correlation to lagged changes in CAN show more statistical significance.
The pseudo R\textsuperscript{2} calculated for the QVAR indicate a low
goodness of fit for both CAN and wheat prices, showing that little
variation in price changes are explained in the model. In general,
effects vary across quantile showing stronger effects in the median
regression, dissipating in the tails.

\subsection{Joint price distribution}\label{joint-price-distribution}

The estimations of the marginal distribution are used to evaluate the
joint distribution. We estimate \(F_{N}\) and \(F_{w}\) by first
deriving the quantiles that provide the best fit for the estimation. The
obtained quantile distributions are used to estimate the copula
connecting the marginal distributions. The copula
\(C\left( F_{w} \middle| F_{N} \right)\) indicates the contemporaneous
co-dependence between the marginal distributions and is estimated using
a quantile regression on the inverse marginal distributions using a
conditional distribution function \((D)\) of the form:

\[\begin{array}{r}
D\left( q_{w} \middle| q_{N} \right) = \mathbf{\beta}_{\mathbf{0}}\left( q_{w} \right) + \mathbf{\beta}_{\mathbf{1}}\left( q_{w} \right)q_{N}\#(25)
\end{array}\]

\protect\phantomsection\label{_Ref214957457}{}\emph{\hfill\break
}

\emph{Table 6.} \emph{Quantile estimation of conditional copula}

{\def\LTcaptype{none} % do not increment counter
\begin{longtable}[]{@{}
  >{\raggedright\arraybackslash}p{(\linewidth - 10\tabcolsep) * \real{0.1667}}
  >{\raggedright\arraybackslash}p{(\linewidth - 10\tabcolsep) * \real{0.1667}}
  >{\raggedright\arraybackslash}p{(\linewidth - 10\tabcolsep) * \real{0.1667}}
  >{\raggedright\arraybackslash}p{(\linewidth - 10\tabcolsep) * \real{0.1667}}
  >{\raggedright\arraybackslash}p{(\linewidth - 10\tabcolsep) * \real{0.1667}}
  >{\raggedright\arraybackslash}p{(\linewidth - 10\tabcolsep) * \real{0.1667}}@{}}
\toprule\noalign{}
\begin{minipage}[b]{\linewidth}\raggedright
\end{minipage} & \begin{minipage}[b]{\linewidth}\raggedright
\end{minipage} & \begin{minipage}[b]{\linewidth}\raggedright
\end{minipage} & \begin{minipage}[b]{\linewidth}\raggedright
\[q_{w}\]
\end{minipage} & \begin{minipage}[b]{\linewidth}\raggedright
\end{minipage} & \begin{minipage}[b]{\linewidth}\raggedright
\end{minipage} \\
\midrule\noalign{}
\endhead
\bottomrule\noalign{}
\endlastfoot
& 0.1 & 0.3 & 0.5 & 0.7 & 0.9 \\
Intercept & 0.036 & 0.224*** & 0.520*** & 0.690*** & 0.868*** \\
& (0.024) & (0.042) & (0.063) & (0.044) & (0.037) \\
\(q_{N}\) & 0.148*** & 0.129** & -0.034 & 0.000 & 0.053 \\
& (0.054) & (0.073) & (0.093) & (0.078) & (0.075) \\
\end{longtable}
}

Note: \(q_{w}\) and \(q_{N}\) denote the quantile distribution of wheat
and CAN. Standard Errors in parenthesis are calculated by bootstrapping;
stars indicate significance at * = 0.1, ** = 0.05, and *** = 0.01
levels.

Results displayed in \hyperref[_Ref214957457]{Table 6} show significant
positive coefficients in the lower quantiles (0.1 and 0.3), indicating a
positive contemporaneous codependence of CAN and wheat price changes.
For the median and higher quantiles we find no evidence against
independence.

\subsection{Contract pricing}\label{contract-pricing}

We simulate 1.000 price paths of 24 time periods (one year). Simulations
for CAN and wheat prices based on the estimated models (see also Fig. A1
for density distributions) are used for the contract pricing and profits
calculations. Results from the burn analysis at different strike prices
are presented in \hyperref[_Ref215049461]{Table 7}. We select the strike
prices based on distribution quantiles of simulated CAN prices.
Diminishing premiums at higher strike prices reflecting lower
probability of exercise and expected payoff.
\hyperref[_Ref215737709]{Figure 3} reports the price density
distribution of CAN under the contract conditions across different
strike prices comparing them to the simulated CAN prices.

\protect\phantomsection\label{_Ref215049461}{}\emph{Table 7.}
\emph{Contract pricing by burn analysis at different strike prices}

{\def\LTcaptype{none} % do not increment counter
\begin{longtable}[]{@{}
  >{\raggedright\arraybackslash}p{(\linewidth - 4\tabcolsep) * \real{0.1029}}
  >{\centering\arraybackslash}p{(\linewidth - 4\tabcolsep) * \real{0.1897}}
  >{\centering\arraybackslash}p{(\linewidth - 4\tabcolsep) * \real{0.1677}}@{}}
\toprule\noalign{}
\begin{minipage}[b]{\linewidth}\raggedright
\textbf{Quantile}
\end{minipage} & \begin{minipage}[b]{\linewidth}\centering
\textbf{Strike price (€/kg)}
\end{minipage} & \begin{minipage}[b]{\linewidth}\centering
\textbf{Premium (€/kg)}
\end{minipage} \\
\midrule\noalign{}
\endhead
\bottomrule\noalign{}
\endlastfoot
0.5 & 0.355 & 0.011 \\
0.6 & 0.361 & 0.009 \\
0.7 & 0.366 & 0.007 \\
0.8 & 0.375 & 0.004 \\
0.9 & 0.389 & 0.002 \\
\end{longtable}
}

\includegraphics[width=6.5in,height=4.97153in]{media/image4.png}

\protect\phantomsection\label{_Ref215737709}{}\emph{Figure 3. Price
density distributions of simulated CAN prices and CAN prices paid with
the contract at different strike levels}

\subsection{\texorpdfstring{Profits and Utility levels
}{Profits and Utility levels }}\label{profits-and-utility-levels}

Based on our decision framework, we assume that the decision-maker can
adopt a forward contract at the start of the season, thus fixing the CAN
price, or adopt the adjusted contract with different price caps.
Alternatively, CAN is bought in early march at market price (no
contract). Simulated profits for the different strategies are calculated
under broad and narrow framings (see Annex 3 for summary statistics).
Results are displayed in \hyperref[_Ref216184323]{Table 8} and
\hyperref[_Ref216184771]{Table 9} for EUT and CPT preferences under
broad framing, and in \hyperref[_Ref216184995]{Table 10} and
\hyperref[_Ref216185008]{Table 11}. EUT results are reported as
Certainty Equivalents (CE) in Euro/ha across different values of
absolute risk aversion \((\delta)\). CPT results are reported as final
prospect values \((V)\). Alternative parameter levels for risk aversion
\((a)\), loss aversion \((\lambda)\) and probability decision weights
\((\gamma)\) are included. The choice of the reference point is
different for broad and narrow framing. When the contract is considered
as a risk management tool, we adopt expected outcome as a reference
point. When it is considered as a stand-alone investment, the reference
point is zero.

Under broad framing, the producer evaluates each contract in the context
of the entire farm profits. Results for the EUT framework suggest that
across all levels of risk aversion (except for extremely risk averse
levels), the adjusted contract delivers the highest CE, indicating the
viability of this strategy when framed as a risk management tool. In
general, however, the differences in CEs are small across the different
strategies. Interestingly, the forward contract yields the lowest CE
across all the \textbackslash delta values, indicating that, under broad
framing, this would be the least preferred strategy. This is explained
by the fact that by fixing input prices, the producer is losing the
natural price hedge (i.e. given the correlation between fertilizers and
wheat prices, when prices are high farmers get higher output prices and
pay higher input prices, the contrary happens when prices are low). This
causes total profits to fluctuate more strongly.

Results for the broad framing using a CPT framework display negative
values for all the strategies, meaning that the prospects are perceived
as a loss relative to the expected profits under the forward contract.
This is consistent across alternative parametrization with the only
exception for the atypical case of loss-seeking behavior
\((\lambda = 0.5)\). This suggests that the adjusted contract is not
attractive to farmers when considering CPT preferences, in contrast to
EUT results.

When the contracts are narrowly framed as risky investments, producers
consider only the contract's payoff and cost, with the reference point
becoming zero. Here, since it is framed as an investment decision, the
``no contract'' strategy is not analyzed. Results for the EUT framework
indicate that forward pricing yield the highest CE values across the
different risk aversion levels. This result stands in direct contrast to
the broad framing scenario. Under narrow framing, profit variability is
driven solely by fertilizer price risk. In contrast, under broad
framing, the correlation between fertilizer and wheat prices makes the
adoption of a price-fixing strategy unattractive to farmers, as it
eliminates the natural price hedge. Consequently, while forward
contracts are disfavored under broad framing, they are preferred under
narrow framing. Similar to the broad framing case, results for the CPT
framework display negative values. Overall, both EUT and CPT results
under narrow framing suggest that the adjusted contract is not a viable
option and might face low adoption.

\protect\phantomsection\label{_Ref216184323}{}\emph{Table 8. Certainty
Equivalents (CE) for different level of absolute risk aversion}
\((\delta)\) \emph{under broad framing.}

{\def\LTcaptype{none} % do not increment counter
\begin{longtable}[]{@{}
  >{\raggedright\arraybackslash}p{(\linewidth - 14\tabcolsep) * \real{0.1094}}
  >{\centering\arraybackslash}p{(\linewidth - 14\tabcolsep) * \real{0.0984}}
  >{\centering\arraybackslash}p{(\linewidth - 14\tabcolsep) * \real{0.1376}}
  >{\centering\arraybackslash}p{(\linewidth - 14\tabcolsep) * \real{0.1274}}
  >{\centering\arraybackslash}p{(\linewidth - 14\tabcolsep) * \real{0.1274}}
  >{\centering\arraybackslash}p{(\linewidth - 14\tabcolsep) * \real{0.1520}}
  >{\centering\arraybackslash}p{(\linewidth - 14\tabcolsep) * \real{0.1278}}
  >{\centering\arraybackslash}p{(\linewidth - 14\tabcolsep) * \real{0.1200}}@{}}
\toprule\noalign{}
\begin{minipage}[b]{\linewidth}\raggedright
\end{minipage} & \begin{minipage}[b]{\linewidth}\centering
\textbf{No contract}
\end{minipage} & \begin{minipage}[b]{\linewidth}\centering
\textbf{Forward contract}
\end{minipage} & \begin{minipage}[b]{\linewidth}\centering
\end{minipage} & \begin{minipage}[b]{\linewidth}\centering
\end{minipage} & \begin{minipage}[b]{\linewidth}\centering
\textbf{Adjusted contract}
\end{minipage} & \begin{minipage}[b]{\linewidth}\centering
\end{minipage} & \begin{minipage}[b]{\linewidth}\centering
\end{minipage} \\
\midrule\noalign{}
\endhead
\bottomrule\noalign{}
\endlastfoot
\textbf{Strike} & \textbf{-} & \textbf{-} & \textbf{0.355} &
\textbf{0.361} & \textbf{0.368} & \textbf{0.376} & \textbf{0.388} \\
\(\delta\mathbf{\ }\)\textbf{=} 0.00001 & 1162.43 & 1160.01 & 1162.31 &
\textbf{1162.57} & 1162.95 & 1162.21 & 1162.41 \\
\(\delta\) = 0.00005 & 1160.47 & 1157.96 & 1160.29 & 1160.56 &
\textbf{1160.95} & 1160.22 & 1160.43 \\
\(\delta\) = 0.0001 & 1158.02 & 1155.38 & 1157.77 & 1158.05 &
\textbf{1158.45} & 1157.73 & 1157.95 \\
\(\delta\) = 0.0005 & 1138.43 & 1134.79 & 1137.57 & 1137.94 &
\textbf{1138.44} & 1137.80 & 1138.14 \\
\(\delta\) = 0.001 & \textbf{1113.95} & 1109.05 & 1112.32 & 1112.81 &
1113.42 & 1112.90 & 1113.37 \\
\end{longtable}
}

Note: CE in EUR/ha; the highest CE for each risk aversion level is
highlighted in bold.

\protect\phantomsection\label{_Ref216184771}{}\emph{Table 9.} \emph{CPT
results under broad framing}

{\def\LTcaptype{none} % do not increment counter
\begin{longtable}[]{@{}
  >{\raggedright\arraybackslash}p{(\linewidth - 14\tabcolsep) * \real{0.2463}}
  >{\centering\arraybackslash}p{(\linewidth - 14\tabcolsep) * \real{0.0962}}
  >{\centering\arraybackslash}p{(\linewidth - 14\tabcolsep) * \real{0.1364}}
  >{\centering\arraybackslash}p{(\linewidth - 14\tabcolsep) * \real{0.0958}}
  >{\centering\arraybackslash}p{(\linewidth - 14\tabcolsep) * \real{0.0958}}
  >{\centering\arraybackslash}p{(\linewidth - 14\tabcolsep) * \real{0.1382}}
  >{\centering\arraybackslash}p{(\linewidth - 14\tabcolsep) * \real{0.0958}}
  >{\centering\arraybackslash}p{(\linewidth - 14\tabcolsep) * \real{0.0954}}@{}}
\toprule\noalign{}
\begin{minipage}[b]{\linewidth}\raggedright
\end{minipage} & \begin{minipage}[b]{\linewidth}\centering
\textbf{No contract}
\end{minipage} & \begin{minipage}[b]{\linewidth}\centering
\textbf{Forward contract}
\end{minipage} & \begin{minipage}[b]{\linewidth}\centering
\end{minipage} & \begin{minipage}[b]{\linewidth}\centering
\end{minipage} & \begin{minipage}[b]{\linewidth}\centering
\textbf{Adjusted contract}
\end{minipage} & \begin{minipage}[b]{\linewidth}\centering
\end{minipage} & \begin{minipage}[b]{\linewidth}\centering
\end{minipage} \\
\midrule\noalign{}
\endhead
\bottomrule\noalign{}
\endlastfoot
\textbf{Strike} & \textbf{-} & \textbf{-} & \textbf{0.355} &
\textbf{0.361} & \textbf{0.368} & \textbf{0.376} & \textbf{0.388} \\
\(a = 0.334;\ \lambda = 1.574;\ \gamma = 0.571\ \ \) & -2.43 & -2.48 &
-2.44 & -2.42 & -2.44 & -2.46 & -2.45 \\
\(a = 0.1\) & -0.47 & -0.48 & -0.48 & -0.48 & -0.48 & -0.48 & -0.48 \\
\(a = 0.5\) & -7.18 & -7.37 & -7.21 & -7.16 & -7.19 & -7.23 & -7.23 \\
\(\lambda = 0.5\) & 0.93 & 0.92 & 0.93 & 0.95 & 0.93 & 0.91 & 0.91 \\
\(\lambda = 2\) & -3.76 & -3.83 & -3.78 & -3.76 & -3.78 & -3.79 &
-3.79 \\
\(\gamma = 0.4\ \ \) & -2.86 & -2.92 & -2.89 & -2.86 & -2.89 & -2.91 &
-2.91 \\
\(\gamma = 0.7\ \ \) & -2.25 & -2.30 & -2.26 & -2.25 & -2.26 & -2.27 &
-2.27 \\
\end{longtable}
}

Note: the assumed reference point is the forward contract expected
outcome.

\protect\phantomsection\label{_Ref216184995}{}\emph{Table 10.}
\emph{Certainty Equivalents (CE) for different level of absolute risk
aversion} \((\delta)\) \emph{under narrow framing.}

{\def\LTcaptype{none} % do not increment counter
\begin{longtable}[]{@{}
  >{\raggedright\arraybackslash}p{(\linewidth - 12\tabcolsep) * \real{0.0945}}
  >{\centering\arraybackslash}p{(\linewidth - 12\tabcolsep) * \real{0.1406}}
  >{\centering\arraybackslash}p{(\linewidth - 12\tabcolsep) * \real{0.1530}}
  >{\centering\arraybackslash}p{(\linewidth - 12\tabcolsep) * \real{0.1530}}
  >{\centering\arraybackslash}p{(\linewidth - 12\tabcolsep) * \real{0.1530}}
  >{\centering\arraybackslash}p{(\linewidth - 12\tabcolsep) * \real{0.1530}}
  >{\centering\arraybackslash}p{(\linewidth - 12\tabcolsep) * \real{0.1530}}@{}}
\toprule\noalign{}
\begin{minipage}[b]{\linewidth}\raggedright
\end{minipage} & \begin{minipage}[b]{\linewidth}\centering
\textbf{Forward contract}
\end{minipage} & \begin{minipage}[b]{\linewidth}\centering
\end{minipage} & \begin{minipage}[b]{\linewidth}\centering
\end{minipage} & \begin{minipage}[b]{\linewidth}\centering
\textbf{Adjusted contract}
\end{minipage} & \begin{minipage}[b]{\linewidth}\centering
\end{minipage} & \begin{minipage}[b]{\linewidth}\centering
\end{minipage} \\
\midrule\noalign{}
\endhead
\bottomrule\noalign{}
\endlastfoot
\textbf{Strike} & - & \textbf{0.355} & \textbf{0.361} & \textbf{0.368} &
\textbf{0.376} & \textbf{0.388} \\
\(\delta\mathbf{\ }\)\textbf{=} 0.00001 & \textbf{2.39} & -0.1 & 0.15 &
0.53 & -0.21 & -0.02 \\
\(\delta\) = 0.00005 & \textbf{2.38} & -0.11 & 0.15 & 0.53 & -0.21 &
-0.02 \\
\(\delta\) = 0.0001 & \textbf{2.37} & -0.11 & 0.15 & 0.53 & -0.22 &
-0.02 \\
\(\delta\) = 0.0005 & \textbf{2.29} & -0.15 & 0.12 & 0.51 & -0.23 &
-0.03 \\
\(\delta\) = 0.001 & \textbf{2.18} & -0.19 & 0.08 & 0.48 & -0.25 &
-0.04 \\
\end{longtable}
}

Note: CE in Euro/ha; the highest CE for each level of risk aversion is
highlighted in bold.

\protect\phantomsection\label{_Ref216185008}{}\emph{Table 11.} \emph{CPT
results under narrow framing}

{\def\LTcaptype{none} % do not increment counter
\begin{longtable}[]{@{}
  >{\raggedright\arraybackslash}p{(\linewidth - 12\tabcolsep) * \real{0.2516}}
  >{\centering\arraybackslash}p{(\linewidth - 12\tabcolsep) * \real{0.1422}}
  >{\centering\arraybackslash}p{(\linewidth - 12\tabcolsep) * \real{0.1212}}
  >{\centering\arraybackslash}p{(\linewidth - 12\tabcolsep) * \real{0.1212}}
  >{\centering\arraybackslash}p{(\linewidth - 12\tabcolsep) * \real{0.1212}}
  >{\centering\arraybackslash}p{(\linewidth - 12\tabcolsep) * \real{0.1212}}
  >{\centering\arraybackslash}p{(\linewidth - 12\tabcolsep) * \real{0.1213}}@{}}
\toprule\noalign{}
\begin{minipage}[b]{\linewidth}\raggedright
\end{minipage} & \begin{minipage}[b]{\linewidth}\centering
\textbf{Forward contract}
\end{minipage} &
\multicolumn{5}{>{\centering\arraybackslash}p{(\linewidth - 12\tabcolsep) * \real{0.6062} + 8\tabcolsep}@{}}{%
\begin{minipage}[b]{\linewidth}\centering
\textbf{Adjusted contract}
\end{minipage}} \\
\midrule\noalign{}
\endhead
\bottomrule\noalign{}
\endlastfoot
\textbf{Strike} & \textbf{-} & \textbf{0.355} & \textbf{0.361} &
\textbf{0.368} & \textbf{0.376} & \textbf{0.388} \\
\(a = 0.334;\ \lambda = 1.574;\ \gamma = 0.571\ \ \) & \textbf{-0.74} &
-0.88 & -0.90 & -0.85 & -1.03 & -0.99 \\
\(a = 0.1\) & \textbf{-0.15} & -0.53 & -0.60 & -0.69 & -0.86 & -1.01 \\
\(a = 0.5\) & -1.72 & -1.21 & -1.13 & -0.93 & -1.11 & \textbf{-0.91} \\
\(\lambda = 0.5\) & \textbf{0.52} & 0.26 & 0.19 & 0.16 & 0.03 & -0.05 \\
\(\lambda = 2\) & \textbf{-1.24} & -1.33 & -1.33 & -1.25 & -1.45 &
-1.37 \\
\(\gamma = 0.4\ \ \) & -1.15 & -0.79 & -0.79 & -0.68 & -0.81 &
\textbf{-0.76} \\
\(\gamma = 0.7\ \ \) & \textbf{-0.51} & -0.98 & -1.01 & -0.99 & -1.2 &
-1.15 \\
\end{longtable}
}

Note: the assumed reference point is zero.

\section{Discussion}\label{discussion}

Our results highlight meaningful differences in how producers evaluate
input price risk-management tools, particularly when considered through
the different lenses of broad and narrow framing. Under broad framing,
the contract is assessed as part of the producer's overall farm
portfolio, leading to choices that reflect improvements in total
profits. In contrast, narrow framing focuses on the contract in
isolation, considering only its standalone payoff and cost structure.
The observed variation in results outcomes suggests that producers'
contract decisions depend not only on the underlying decision theory but
also on the cognitive frame through which the decision is made, as
suggested by findings on crop insurance (Babcock, 2015; Dalhaus et al.,
2020). Except for limited evidence of EUT values under broad framing,
the results suggests that, under both framing, the adjusted contract
leads to no meaningful improvement for input price risk. A more holistic
approach that considers both input and output price risk might be better
suited.

The results from the Quantile Vector Autoregression (QVAR) combined with
a copula allowed us to investigate the structure of the multivariate
distribution of Calcium Ammonium Nitrate (CAN) and wheat prices. We find
positive correlations between price changes and changes in their own and
cross lagged changes, especially around the median, with the correlation
fading for lower and higher quantiles of the price distribution. This
approach allowed us to consider the dependency of the two prices on each
other in the forecast model. Nonetheless, simulated prices for CAN
present a shifted price distribution (Fig. A1), implying that our
forecasted prices present higher average values than the observed
empirical distribution. This might bias our quantitative results, but
since we base both the contract pricing and the profit calculation on
forecasted prices, the qualitative interpretation does not lose
validity.

The stochastic simulation incorporates also production risk through the
estimation of a stochastic production function. We use data from the
EPIC-IIASA crop model that simulates yield levels for specific spatial
coordinates, that allows the site specific modelling employed here.
However, a drawback of the data we use is that the crop model
underestimate average winter wheat yields and their standard deviation
(Balkovič et al., 2013). This implies that we might underestimate total
profits and production risk in the broad framing calculations. It is
important to note also that only yields from an optimal nitrogen
application range are simulated. This implies that our stochastic
production function estimation does not capture diminishing marginal
productivity of fertilizers. Nonetheless, we assume that our
representative farm applies an optimal N fertilizer level, mitigating
the implication of the biased estimation. Our methodology is able to
model price dependencies, but does not capture wheat yield and wheat
price correlation (i.e. the level of ``natural hedge''). This omission
is unlikely to materially affect results, given our focus on a small
production region where local shocks do not exert measurable feedback
effects on price formation.

Lastly a note on our modelling approach, decision theories models, and
parameters specifications is needed. First, we restricted our analysis
on a single agricultural activity (winter wheat production) and assumed
no other costs beside input prices. This simplification was employed for
analysis purposes but might not reflect the complexity of real word
farming activities. In addition, strategies to reduce dependency on
nitrogen fertilizers and fertilizers saving technologies, can also have
the dual benefit of reducing input price risk and environmental impact
(Krein et al., 2023).

The EUT preferences are evaluated using the Mean-Variance (M-V)
approximation (Eq. 1). This model has been widely used in the literature
for its simplicity and easy interpretation, especially for hedging
decisions (Spada et al., 2025). Nonetheless, this approach has been
criticized as imposes restrictive behavioral assumption such as the
independence to wealth levels, and the reliance on only the first two
moment of the outcome distribution. Moreover, the selected coefficients
may not fully cover the range of farmers attitudes, affecting in turn
the interpretation of our results. Thus, while this framework
facilitates the tractability, it may not fully capture the heterogeneity
and wealth-dependent nature of actual farmers behavior.

A similar argument can be made for CPT preferences. First, final
prospect values heavily depends on the model specification and
references point. Further, heterogeneity and instability of the
parameters also affect results (Finger et al., 2024). The parameters
employed in the analysis have been taken from empirical evidence on
German farmers (Rommel et al., 2023) and robustness to alternative
parametrization has been explored, indicating similar results. It is
worth noting anyway that our point estimates might not reflect well the
heterogeneity in elicited parameter that farmers display. In addition,
externally imposed reference points have been frequently criticized
(Kőszegi \& Rabin, 2006, 2007). Producer reference points are
unobservable, they can be highly heterogeneous among farmers and shift
in time. Consequently, the results must be interpreted with caution
given the sensitivity to the selected parameters and reference points.

\section{Conclusions}\label{conclusions}

In summary, we propose a maximum-price adjustment to a forward contract
for input price risk management. This adjustment can better align to
producers\textquotesingle{} actual behavioral preferences, thus making
it more attractive to farmers. We estimate both production risk, using a
stochastic production function, and price risk, by providing evidence on
the quantile correlation between fertilizer and wheat prices. We use the
estimated risks in a Montecarlo simulation to assess the feasibility of
the introduction of the adjusted forward contract for input price risk
management. This approach enables us to simulate price and production
risk for utility levels calculations under different decision theories
and framing scenarios.

\section{We evaluate the adjusted contracts using Mean-Variance values
under Expected Utility Theory and prospect values under Cumulative
Prospect Theory. We also evaluate the outcomes under different decision
framing: producers might frame the adoption of the contract in the
context of a whole farm decision, considering it as a risk management
tool (broad framing); alternatively, producers might narrowly frame the
adoption as a standalone investment decision, considering it, in fact,
as an investment. Results from the Montecarlo simulation, that
incorporates both production and price risk, indicate that the contract
delivers limited utility gains for farmers under broad framing, and that
is not attractive when considered under narrow framing, suggesting that
its real-word attractiveness may be low and that adoption could be
limited if the contract were
introduced.References}\label{we-evaluate-the-adjusted-contracts-using-mean-variance-values-under-expected-utility-theory-and-prospect-values-under-cumulative-prospect-theory.-we-also-evaluate-the-outcomes-under-different-decision-framing-producers-might-frame-the-adoption-of-the-contract-in-the-context-of-a-whole-farm-decision-considering-it-as-a-risk-management-tool-broad-framing-alternatively-producers-might-narrowly-frame-the-adoption-as-a-standalone-investment-decision-considering-it-in-fact-as-an-investment.-results-from-the-montecarlo-simulation-that-incorporates-both-production-and-price-risk-indicate-that-the-contract-delivers-limited-utility-gains-for-farmers-under-broad-framing-and-that-is-not-attractive-when-considered-under-narrow-framing-suggesting-that-its-real-word-attractiveness-may-be-low-and-that-adoption-could-be-limited-if-the-contract-were-introduced.references}

Babcock, B. A. (2015). Using cumulative prospect theory to explain
anomalous crop insurance coverage choice. \emph{American Journal of
Agricultural Economics}, \emph{97}(5), 1371--1384.
https://doi.org/10.1093/ajae/aav032

Balkovič, J., van der Velde, M., Schmid, E., Skalský, R., Khabarov, N.,
Obersteiner, M., Stürmer, B., \& Xiong, W. (2013). Pan-European crop
modelling with EPIC: Implementation, up-scaling and regional crop yield
validation. \emph{Agricultural Systems}, \emph{120}, 61--75.
https://doi.org/10.1016/j.agsy.2013.05.008

Barberis, N., Huang, M., \& Thaler, R. H. (2006). Individual
Preferences, Monetary Gambles, and Stock Market Participation: A Case
for Narrow Framing. \emph{American Economic Review}, \emph{96}(4),
1069--1090. https://doi.org/10.1257/aer.96.4.1069

Bocquého, G., Jacquet, F., \& Reynaud, A. (2014). Expected utility or
prospect theory maximisers? Assessing farmers' risk behaviour from
field-experiment data. \emph{European Review of Agricultural Economics},
\emph{41}(1), 135--172. https://doi.org/10.1093/erae/jbt006

Bougherara, D., Gassmann, X., Piet, L., \& Reynaud, A. (2017).
Structural estimation of farmers' risk and ambiguity preferences: A
field experiment. \emph{European Review of Agricultural Economics},
\emph{44}(5), 782--808. https://doi.org/10.1093/erae/jbx011

Couture, S., Lemarié, S., Teyssier, S., \& Toquebeuf, P. (2024). The
value of information under ambiguity: A theoretical and experimental
study on pest management in agriculture. \emph{Theory and Decision},
\emph{96}(1), 19--47. https://doi.org/10.1007/s11238-023-09942-y

Dalhaus, T., Barnett, B. J., \& Finger, R. (2020). Behavioral weather
insurance: Applying cumulative prospect theory to agricultural insurance
design under narrow framing. \emph{PLOS ONE}, \emph{15}(5), e0232267.
https://doi.org/10.1371/journal.pone.0232267

Dong, F., Beckman, J., \& Arita, S. (2025). Unraveling demand and supply
shocks in the U.S. nitrogen market. \emph{Applied Economic Perspectives
and Policy}, \emph{n/a}(n/a). https://doi.org/10.1002/aepp.70004

European Commission. (2022). \emph{Food security: The Commission
addresses the availability and affordability of fertilisers in the EU
and globally} {[}Press release{]}.
https://ec.europa.eu/commission/presscorner/detail/en/ip\_22\_6564

FAO. (2024). \emph{Inorganic fertilizers---2002-2022} (FAOSTAT
Analytical Brief No. 90). https://doi.org/10.4060/cd1485en

Finger, R. (2012). Nitrogen use and the effects of nitrogen taxation
under consideration of production and price risks. \emph{Agricultural
Systems}, \emph{107}, 13--20. https://doi.org/10.1016/j.agsy.2011.12.001

Finger, R. (2013). Investigating the performance of different estimation
techniques for crop yield data analysis in crop insurance applications.
\emph{Agricultural Economics}, \emph{44}(2), 217--230.
https://doi.org/10.1111/agec.12005

Finger, R., Garcia, V., McCallum, C., \& Rommel, J. (2024). A note on
European farmers' preferences under cumulative prospect theory.
\emph{Journal of Agricultural Economics}, \emph{75}(1), 465--472.
https://doi.org/10.1111/1477-9552.12565

Finger, R., \& Hediger, W. (2008). The Application of Robust Regression
to a Production Function Comparison. \emph{The Open Agriculture
Journal}, \emph{2}(1). https://doi.org/10.2174/1874331500802010090

Finger, R., \& Schmid, S. (2008). Modeling agricultural production risk
and the adaptation to climate change. \emph{Agricultural Finance
Review}, \emph{68}(1), 25--41. https://doi.org/10.1108/00214660880001217

Ghirardato, P., Maccheroni, F., \& Marinacci, M. (2004). Differentiating
ambiguity and ambiguity attitude. \emph{Journal of Economic Theory},
\emph{118}(2), 133--173. https://doi.org/10.1016/j.jet.2003.12.004

Gilboa, I., \& Schmeidler, D. (1989). Maxmin expected utility with
non-unique prior. \emph{Journal of Mathematical Economics},
\emph{18}(2), 141--153. https://doi.org/10.1016/0304-4068(89)90018-9

Hardaker, J. B., Huirne, R. B. M., Anderson, J. R., \& Lien, G. (2015).
\emph{Coping With Risk in Agriculture} (3rd ed.). CABI Publishing.
https://www.cabidigitallibrary.org/doi/book/10.1079/9781780645742.0000

Hernández-Ochoa, I. M., Gaiser, T., Hüging, H., \& Ewert, F. (2023).
Yield components and yield quality of old and modern wheat cultivars as
affected by cultivar release date, N fertilization and environment in
Germany. \emph{Field Crops Research}, \emph{302}, 109094.
https://doi.org/10.1016/j.fcr.2023.109094

IFA. (2022). \emph{Fertilizer use by crop and country for the 2017-2018
period}. International Fertilizer Association (IFA).
https://www.ifastat.org/consumption/fertilizer-use-by-crop

Just, R. E., \& Pope, R. D. (1978). Stochastic specification of
production functions and economic implications. \emph{Journal of
Econometrics}, \emph{7}(1), 67--86.
https://doi.org/10.1016/0304-4076(78)90006-4

Just, R. E., \& Pope, R. D. (1979). Production Function Estimation and
Related Risk Considerations. \emph{American Journal of Agricultural
Economics}, \emph{61}(2), 276--284. https://doi.org/10.2307/1239732

Knight, F. H. (1921). \emph{Risk, uncertainty and profit} (Vol. 31).
Houghton Mifflin.

Koenker, R. (2005). \emph{Quantile Regression}. Cambridge University
Press. https://doi.org/10.1017/CBO9780511754098

Koenker, R., \& Machado, J. A. F. (1999). Goodness of Fit and Related
Inference Processes for Quantile Regression. \emph{Journal of the
American Statistical Association}, \emph{94}(448), 1296--1310.
https://doi.org/10.1080/01621459.1999.10473882

Koenker, R., \& Xiao, Z. (2006). Quantile Autoregression. \emph{Journal
of the American Statistical Association}, \emph{101}(475), 980--990.
https://doi.org/10.1198/016214506000000672

Kőszegi, B., \& Rabin, M. (2007). Reference-Dependent Risk Attitudes.
\emph{American Economic Review}, \emph{97}(4), 1047--1073.
https://doi.org/10.1257/aer.97.4.1047

Krein, D. D. C., Rosseto, M., Cemin, F., Massuda, L. A., \& Dettmer, A.
(2023). Recent trends and technologies for reduced environmental impacts
of fertilizers: A review. \emph{International Journal of Environmental
Science and Technology}, \emph{20}(11), 12903--12918.
https://doi.org/10.1007/s13762-023-04929-2

Li, J., \& Chavas, J. P. (2023). A dynamic analysis of the distribution
of commodity futures and spot prices. \emph{American Journal of
Agricultural Economics}, \emph{105}(1), 122--143.
https://doi.org/10.1111/ajae.12309

Mattos, F., Garcia, P., \& Pennings, J. M. E. (2008). Probability
weighting and loss aversion in futures hedging. \emph{Journal of
Financial Markets}, \emph{11}(4), 433--452.
https://doi.org/10.1016/j.finmar.2008.04.002

Moritz, L., Spada, R., Rommel, J., Dalhaus, T., \& Cerroni, S. (2026).
Risk preferences and other (ignored) behavioral factors in fertilizer
management decisions: A systematic literature review. \emph{Journal of
Behavioral and Experimental Economics}, \emph{121}, 102524.
https://doi.org/10.1016/j.socec.2026.102524

Nelsen. (2006). \emph{An Introduction to Copulas}. Springer.
https://doi.org/10.1007/0-387-28678-0

O'Neill, B. C., Kriegler, E., Riahi, K., Ebi, K. L., Hallegatte, S.,
Carter, T. R., Mathur, R., \& van Vuuren, D. P. (2014). A new scenario
framework for climate change research: The concept of shared
socioeconomic pathways. \emph{Climatic Change}, \emph{122}(3), 387--400.
https://doi.org/10.1007/s10584-013-0905-2

Prelec, D. (1998). The Probability Weighting Function.
\emph{Econometrica}, \emph{66}(3), 497--527.
https://doi.org/10.2307/2998573

Rommel, J., Sagebiel, J., Baaken, M. C., Barreiro-Hurlé, J., Bougherara,
D., Cembalo, L., Cerjak, M., Čop, T., Czajkowski, M., Espinosa-Goded,
M., Höhler, J., Kuhfuss, L., Lagerkvist, C.-J., Lapierre, M., Lefebvre,
M., Matzdorf, B., Ott, E., Paparella, A., Quendler, E., \ldots{}
Zagórska, K. (2023). Farmers' risk preferences in 11 European farming
systems: A multi-country replication of Bocquého et al. ().
\emph{Applied Economic Perspectives and Policy}, \emph{45}(3),
1374--1399. https://doi.org/10.1002/aepp.13330

Sklar, M. (1959). \emph{Fonctions de Répartition À N Dimensions Et Leurs
Marges}. Université Paris 8.
https://books.google.nl/books?id=nreSmAEACAAJ

Spada, R., Moritz, L., Rommel, J., Cerroni, S., Meuwissen, M. P. M., \&
Dalhaus, T. (2025). \emph{Behavioral Preferences Affecting the Adoption
of Financial Price Risk Management Tools in Agriculture: A Systematic
Literature Review}. https://doi.org/10.22004/ag.econ.356632

Tversky, A., \& Kahneman, D. (1981). The Framing of Decisions and the
Psychology of Choice. \emph{Science}, \emph{211}(4481), 453--458.
https://doi.org/10.1126/science.7455683

Tversky, A., \& Kahneman, D. (1992). Advances in prospect theory:
Cumulative representation of uncertainty. \emph{Journal of Risk and
Uncertainty}, \emph{5}(4), 297--323. https://doi.org/10.1007/BF00122574

Yohai, V. J. (1987). High Breakdown-Point and High Efficiency Robust
Estimates for Regression. \emph{The Annals of Statistics}, \emph{15}(2),
642--656. https://doi.org/10.1214/aos/1176350366

\section{Appendix}\label{appendix}

\textbf{Table A1} \emph{Descriptive statistics for simulated yield data}

{\def\LTcaptype{none} % do not increment counter
\begin{longtable}[]{@{}
  >{\raggedright\arraybackslash}p{(\linewidth - 8\tabcolsep) * \real{0.3594}}
  >{\centering\arraybackslash}p{(\linewidth - 8\tabcolsep) * \real{0.1835}}
  >{\centering\arraybackslash}p{(\linewidth - 8\tabcolsep) * \real{0.1426}}
  >{\centering\arraybackslash}p{(\linewidth - 8\tabcolsep) * \real{0.1142}}
  >{\centering\arraybackslash}p{(\linewidth - 8\tabcolsep) * \real{0.2002}}@{}}
\toprule\noalign{}
\begin{minipage}[b]{\linewidth}\raggedright
\textbf{variable}
\end{minipage} & \begin{minipage}[b]{\linewidth}\centering
\textbf{mean}
\end{minipage} & \begin{minipage}[b]{\linewidth}\centering
\textbf{sd}
\end{minipage} & \begin{minipage}[b]{\linewidth}\centering
\textbf{min}
\end{minipage} & \begin{minipage}[b]{\linewidth}\centering
\textbf{max}
\end{minipage} \\
\midrule\noalign{}
\endhead
\bottomrule\noalign{}
\endlastfoot
N (kg/ha) & 167.17 & 55.97 & 10 & 250 \\
yield (t/ha) & 7.39 & 2.25 & 1.16 & 14.01 \\
historical weather & 0.53 & - & 0 & 1 \\
cmip6\_ipsl\_ssp126 & 0.12 & - & 0 & 1 \\
cmip6\_ipsl\_ssp585 & 0.12 & - & 0 & 1 \\
cmip6\_mpi\_ssp126 & 0.12 & - & 0 & 1 \\
cmip6\_mpi\_ssp585 & 0.12 & - & 0 & 1 \\
rainfed & 0.86 & - & 0 & 1 \\
irrigated & 0.14 & - & 0 & 1 \\
rot\_mono & 0.33 & - & 0 & 1 \\
rot\_1 & 0.15 & - & 0 & 1 \\
rot\_2 & 0.07 & - & 0 & 1 \\
rot\_3 & 0.07 & - & 0 & 1 \\
rot\_4 & 0.15 & - & 0 & 1 \\
rot\_6 & 0.08 & - & 0 & 1 \\
rot\_7 & 0.15 & - & 0 & 1 \\
til\_bau & 0.57 & - & 0 & 1 \\
til\_min\_cons & 0.1 & - & 0 & 1 \\
til\_contill\_r00 & 0.11 & - & 0 & 1 \\
til\_contill\_r30 & 0.07 & - & 0 & 1 \\
til\_contill\_r60 & 0.09 & - & 0 & 1 \\
til\_contill\_r90 & 0.07 & - & 0 & 1 \\
\end{longtable}
}

Note: outputs from different climate models based on phase 6 of the
Coupled Model Intercomparison Project (cmip6) from the Institut
Pierre-Simone Laplace (ipsl) and Max Planck Institute for Meteorology
(mpi) are considered under different Shared Socioeconomic Pathways
(ssp); ssp126 is a low-emission scenario; ssp585 is a very high emission
scenario (O'Neill et al., 2014). Irrigation is assumed only for N
application rates \(\geq\) 180 kgN/ha; in the case of a monoculture
cropping system (rot\_mono), conventional tillage with different
percentages of crop residue are assumed (til\_contill\_r) across 0, 30,
60, and 90 percent; in the case of different crop rotations (rot\_1-7),
conventional tillage with business-as-usual crop residue (til\_bau) and
minimum tillage with high crop residues (til\_min\_cons) are considered;
yield levels are obtained converting dry-matter yields considering a
water content of 12\%; for more information about the EPIC-IIASA model
specifications and calibration see Balkovič et al. (2013).

\textbf{\hfill\break
}

\textbf{Table A2} \emph{Summary statistics for price data}

{\def\LTcaptype{none} % do not increment counter
\begin{longtable}[]{@{}
  >{\raggedright\arraybackslash}p{(\linewidth - 8\tabcolsep) * \real{0.2000}}
  >{\centering\arraybackslash}p{(\linewidth - 8\tabcolsep) * \real{0.2000}}
  >{\centering\arraybackslash}p{(\linewidth - 8\tabcolsep) * \real{0.2000}}
  >{\centering\arraybackslash}p{(\linewidth - 8\tabcolsep) * \real{0.2000}}
  >{\centering\arraybackslash}p{(\linewidth - 8\tabcolsep) * \real{0.2000}}@{}}
\toprule\noalign{}
\begin{minipage}[b]{\linewidth}\raggedright
\textbf{Statistic}
\end{minipage} & \begin{minipage}[b]{\linewidth}\centering
\textbf{log\_CAN}
\end{minipage} & \begin{minipage}[b]{\linewidth}\centering
\textbf{log\_Wheat}
\end{minipage} & \begin{minipage}[b]{\linewidth}\centering
\textbf{dif\_CAN}
\end{minipage} & \begin{minipage}[b]{\linewidth}\centering
\textbf{dif\_Wheat}
\end{minipage} \\
\midrule\noalign{}
\endhead
\bottomrule\noalign{}
\endlastfoot
Mean & 3.06 & 1.82 & 0 & 0 \\
SD & 1.44 & 0.51 & 0.23 & 0.12 \\
Min & 1.62 & 0.94 & -0.92 & -1.04 \\
Max & 9.29 & 3.8 & 3.03 & 1.08 \\
Skew & 2.52 & 1.25 & 6.75 & 1.15 \\
Kurtosis & 9.3 & 5.51 & 89.6 & 44.23 \\
ADF test & -2.47 & -2.73 & -8.47* & -9.05* \\
KPSS test & 1.45* & 1.42* & 0.07 & 0.08 \\
PP test & -11.19 & -12.8 & -295.74* & -474.91* \\
\end{longtable}
}

Note: statistics rounded to the second decimal of Calcium Ammonium
Nitrate (CAN) and wheat prices (Eur/kg); * indicates the rejection of
the null hypothesis at the 5\% level for the Augmented Dickey-Fuller
(ADF) test, the Kwiatkowski-Phillips-Schmidt-Shin (KPSS), and the
Phillips-Perron (PP) test.

\includegraphics[width=6.5in,height=4.23264in]{media/image5.png}

\textbf{Figure A1.} \emph{Density distributions of simulated and
deflated CAN and wheat prices}

\textbf{Table A3}\emph{. Summary statistics of simulated CAN prices,
costs, and profits}

{\def\LTcaptype{none} % do not increment counter
\begin{longtable}[]{@{}
  >{\raggedright\arraybackslash}p{(\linewidth - 14\tabcolsep) * \real{0.1651}}
  >{\raggedright\arraybackslash}p{(\linewidth - 14\tabcolsep) * \real{0.1807}}
  >{\centering\arraybackslash}p{(\linewidth - 14\tabcolsep) * \real{0.1013}}
  >{\centering\arraybackslash}p{(\linewidth - 14\tabcolsep) * \real{0.0979}}
  >{\centering\arraybackslash}p{(\linewidth - 14\tabcolsep) * \real{0.1230}}
  >{\centering\arraybackslash}p{(\linewidth - 14\tabcolsep) * \real{0.1436}}
  >{\centering\arraybackslash}p{(\linewidth - 14\tabcolsep) * \real{0.0952}}
  >{\centering\arraybackslash}p{(\linewidth - 14\tabcolsep) * \real{0.0934}}@{}}
\toprule\noalign{}
\begin{minipage}[b]{\linewidth}\raggedright
\end{minipage} & \begin{minipage}[b]{\linewidth}\centering
\end{minipage} & \begin{minipage}[b]{\linewidth}\centering
\textbf{strike}
\end{minipage} & \begin{minipage}[b]{\linewidth}\centering
\textbf{mean}
\end{minipage} & \begin{minipage}[b]{\linewidth}\centering
\textbf{median}
\end{minipage} & \begin{minipage}[b]{\linewidth}\centering
\textbf{sd}
\end{minipage} & \begin{minipage}[b]{\linewidth}\centering
\textbf{min}
\end{minipage} & \begin{minipage}[b]{\linewidth}\centering
\textbf{Max}
\end{minipage} \\
\midrule\noalign{}
\endhead
\bottomrule\noalign{}
\endlastfoot
CAN price paid & No contract & - & 0.356 & 0.355 & 0.028 & 0.279 &
0.523 \\
& Forward contract & - & 0.359 & 0.359 & 0 & 0.359 & 0.359 \\
& Adjusted contract & 0.355 & 0.345 & 0.355 & 0.015 & 0.279 & 0.355 \\
& & 0.361 & 0.348 & 0.355 & 0.017 & 0.279 & 0.361 \\
& & 0.368 & 0.35 & 0.355 & 0.019 & 0.279 & 0.368 \\
& & 0.376 & 0.352 & 0.355 & 0.021 & 0.279 & 0.376 \\
& & 0.388 & 0.354 & 0.355 & 0.023 & 0.279 & 0.388 \\
Costs & No contract & - & 263.65 & 262.82 & 20.47 & 206.54 & 387.11 \\
& Forward contract & - & 266.04 & 266.04 & 0 & 266.04 & 266.04 \\
& Adjusted contract & 0.355 & 263.75 & 270.97 & 11.20 & 214.68 &
271.11 \\
& & 0.361 & 263.49 & 268.75 & 12.60 & 212.46 & 273.33 \\
& & 0.368 & 263.12 & 266.52 & 14.18 & 210.24 & 276.29 \\
& & 0.376 & 263.86 & 265.78 & 15.69 & 209.50 & 281.48 \\
& & 0.388 & 263.66 & 264.30 & 17.32 & 208.02 & 288.89 \\
\multirow{7}{=}{Profits broad framing} & No contract & - & 1162.92 &
1122.59 & 312.96 & 279.36 & 2316.85 \\
& Forward contract & - & 1160.53 & 1120.58 & 320.86 & 262.26 &
2395.28 \\
& Adjusted contract & 0.355 & 1162.82 & 1123.57 & 317.80 & 271.21 &
2390.21 \\
& & 0.361 & 1163.07 & 1121.78 & 317.08 & 273.44 & 2387.99 \\
& & 0.368 & 1163.45 & 1121.99 & 316.33 & 275.66 & 2385.02 \\
& & 0.376 & 1162.71 & 1120.81 & 315.62 & 276.4 & 2379.84 \\
& & 0.388 & 1162.90 & 1121.86 & 314.76 & 277.88 & 2372.43 \\
\multirow{7}{=}{Profits narrow framing} & No contract & - & - & - & - &
- & - \\
& Forward contract & - & 2.39 & 3.22 & 20.47 & -121.07 & 59.5 \\
& Adjusted contract & 0.355 & -0.1 & -8.15 & 13.23 & -8.15 & 115.99 \\
& & 0.361 & 0.15 & -5.93 & 11.85 & -5.93 & 113.77 \\
& & 0.368 & 0.53 & -3.7 & 10.31 & -3.7 & 110.81 \\
& & 0.376 & -0.21 & -2.96 & 8.71 & -2.96 & 105.62 \\
& & 0.388 & -0.02 & -1.48 & 6.71 & -1.48 & 98.22 \\
\end{longtable}
}

\end{document}
